\section{A scalar norm tree with inverse weights}

This replacement starts from the frozen fifth-round builder with SHA256
\texttt{68cb3cb1a809766aa102b8c43444d8b38c02760751e773a960489c327e33c12b}.
It preserves the scalar order twenty-eight while removing two rational
coordinates after the full compiler. The point domains, the unnormalized
scalar proxy, the five-contact height theorem, and the inherited global
arithmetic inputs are unchanged. This is a transfer of the scalar
soundness and completeness argument; it is not a bijection between all
old and new scalar witness tuples.

\subsection{The actual source equations}

Retain
\[
 \beta=2401+2754(2X+1)^2,\qquad
 r=\beta(\beta-1)(32\beta+343)>0.
\]
The scalar points are $H=[h]P$ and $Q=S-P=[a-1]P$, with
$h\ne0$ and $a-1\ne0$. Their existing charts, including the rational
inverse of $f(Q)_0$, are retained. In particular the removed identity
output $S=P$ is still excluded; the tangent case $S=-P$ is retained.
Write
\[
 (A,B,C,D_0)=\bigl(f(H)_0,f(H)_1,f(Q)_0,f(Q)_1\bigr),
 \qquad T=\operatorname{Re}(k_Hw).
\]
Here $w=h_S\tau$, $\tau^2=f(Q)$, and $k_H=y(H)/x(H)$ are exactly
the existing fifth-round functions. The letter $D_0$ in this section
denotes a rational scalar coefficient, not the curve coefficient.

Starting with the already required $q_1=r^{-1}$, use three Laurent
steps rather than the former four power coordinates:
\[
 rq_3=q_1^2,\qquad rq_7=q_3^2,\qquad rq_{15}=q_7^2.
\]
Thus $\kappa=q_{15}=r^{-15}$. Introduce the scaled proxy by
\[
 ry=Tq_{15},\qquad\text{so }y=T/r^{16}.
\]
Keep the optimized left subtree, but change the weights in the right
subtree and the final node:
\begin{align*}
 n_0&=1+A^2/r,&n_1&=B^2/r^2+C^2/r,\\
 L&=n_0^2+r n_1^2,&J&=D_0^2+\kappa y^2,\\
 W&=L^2+\kappa J^2.
\end{align*}
The outer scalar pole input remains exactly $b=42rW^2$.
Since $r,\kappa>0$ at the real place, $n_0\ge1$, $L\ge1$ and
$W\ge1$. All the existing absorption of $W^2$ into the pole variables
therefore remains legitimate. No scalar pole equation is changed.

Every displayed operation has a quadratic implementation. For $n_1$
one retains $z=B/r$ and the equation $rn_1=zB+C^2$. The three later
norm nodes each materialize the square of their right argument before
multiplying it by $r$ or $q_{15}$. In particular no cubic product is
silently inlined. The existing quadratic expression for $T$ is retained.

\subsection{Soundness: inverse weights increase the right-branch bounds}

Let $p$ be a consumed prime for $M_{\mathcal S}(\beta)$, where
$\mathcal S$ is the already selected exceptional set. The earlier
three-contact calculation gives
\[
 R=v_p(r)>0\quad\text{and}\quad R\text{ odd},
\]
and $42$ is a $p$-unit. The pole predicate and $b\ne0$ imply
\[
 R+2v_p(W)>0,\qquad v_p(W)>-R/2.
\]
Alternatively, if one starts with only oddness of $R$, the constant
all-left path and the noncancellation calculation below give $v_p(W)\le0$,
so the same pole implication first excludes $R<0$.

For rational inputs, the two summands in
$u^2+r v^2$ have opposite valuation parity, and so do the two summands
in $u^2+r^{-15}v^2$. Consequently the valuation of either norm node is
the minimum of its two summand valuations. This remains valid when one
argument vanishes, with $v_p(0)=+\infty$.

Descending first through $W=L^2+r^{-15}J^2$ gives
\[
 v_p(L)>-R/4,\qquad v_p(J)>29R/4.
\]
Descending through the retained left subtree gives
\[
 v_p(n_0)>-R/8,\qquad v_p(n_1)>-5R/8.
\]
Its two initial nodes then imply
\[
 v_p(A)>7R/16,\qquad
 v_p(B)>11R/16,\qquad
 v_p(C)>3R/16.
\]
The two summands of $n_1=B^2/r^2+C^2/r$ also have opposite valuation
parity because $R$ is odd. The right subtree implies
\[
 v_p(D_0)>29R/8,\qquad v_p(y)>89R/8.
\]
All four coefficients of $f(H)$ and $f(Q)$ therefore have positive
integer valuation. This puts both points in the formal group at every
place above $p$, before any use of the proxy conclusion.

Since $y=T/r^{16}$,
\[
 v_p(T)>\left(16+\frac{89}{8}\right)R=\frac{217}{8}R.
\]
In particular $v_p(T)\ge28$, as $R\ge1$ and the valuation of a
nonzero rational number is integral. The fifth-round two-place proxy
lemma now applies without alteration:
\[
 v_p(T)=v_p\bigl((a-1)/h\bigr).
\]
Its hypotheses have just been obtained from the first four scalar
coefficients. That lemma covers split and inert places, including the
nonunit fixed-coordinate case at 41. Since $h$ is a nonzero integer,
\[
 v_p(a-1)\ge28,
 \qquad M_{\mathcal S}(\beta)^{28}\mid a-1.
\]
Thus the exact order-twenty-eight height consumer is preserved. There
is no new exceptional prime, no use of an unproved exact proxy identity,
and no need to modify any logarithm precision in this scalar argument.

\subsection{Completeness and the order of choices}

Fix an integral input, so $r$ is a fixed positive integer. The inherited
scalar preparation first chooses $H=[h_0]P$ sufficiently deep at the
finite prime support of $42r$, and then chooses $Q=[a-1]P$ sufficiently
deeper. It permits arbitrary finite lower bounds on the coefficients of
$f(H)$ and $f(Q)$ and on $T$: for fixed $H$, as $Q\to O$ one has
$S\to P$, bounded $h_S$, and $w=h_S\tau\to0$, hence $T\to0$.

Here are explicit sufficient new bounds at a prime $p\mid r$, with
$R=v_p(r)\ge1$. Require
\[
 v_p(A),v_p(B),v_p(C)\ge R,\qquad
 v_p(D_0)\ge4R+1,\qquad v_p(y)\ge12R+1.
\]
They give $n_0\equiv1\pmod p$, integral $n_1$, and
$L\equiv1\pmod p$. Moreover
\[
 v_p(J)\ge\min(8R+2,9R+2)=8R+2,
\]
so $v_p(\kappa J^2)\ge R+4>0$, and $W\equiv1\pmod p$.
The bound on $y$ amounts to the finite condition $v_p(T)\ge28R+1$.
It can be imposed after $H$ is fixed. These bounds are compatible with
all prior arbitrary divisibility and size requirements on the indices.

At a prime dividing 42 but not $r$, both $r$ and $\kappa$ are units.
Requiring the four scalar coefficients and $y$ to be divisible by $p$
again gives $W\equiv1\pmod p$. These statements also cover the
ramified and bad primes: the preparation uses sufficiently deep formal
subgroups there, not the good-reduction valuation lemma used for
soundness. At every other prime $42r$ is a unit and $W^2$ contributes
even valuation, as in the original pole proof.

Thus $b=42rW^2$ is positive, nonzero, has square reduced denominator,
and has the same local squareclasses as $42r$ at $2,3,17$. The inherited
scalar preparation and pole theorem supplies the witnesses. With $S$
fixed, all changed leaves are rational functions of $H$ with nonzero
denominators at the prepared point, so the finite lower bounds persist
on a congruence class $h\equiv h_0\pmod L$. The roots and scales are
fixed before the later contact-prime selection, exactly as before.

\subsection{Measured cost, independent checks, and alternatives}

The new power chain uses three coordinates in place of four. The old
coordinate $f(Q)_1/r$ is deleted, since the right subtree now uses
$D_0=f(Q)_1$ directly. All other scalar operations retain their former
coordinate cost. The complete isolated pipeline gives
\[
 \text{raw }413/381,\qquad
 \text{generic compiler }403/371,\qquad
 \boxed{\text{final }399\text{ coordinates and }367\text{ residuals}}.
\]
The two numbers in the first two pairs are coordinates and residuals.
The residuals are quadratic, hence their sum of squares is quartic.
The independent gate replay checks all 352 gate identities and 340
compiler operations. Four further certified eliminations give the
stated final count. These are measurements of the actual implementation,
not estimates from a list of algebraic abbreviations.

The preceding all-inverse-weight prototype also worked, but replacing
the left subtree as well lost one of its cheap quotient realizations.
It compiled to 400/368 and is superseded by the mixed-weight tree.
Simply replacing the old power chain by $r^{-31}$ can increase nominal
precision but saves no coordinate; it does not exploit the valuation
contribution of the two right norm edges. A direct replacement of all
norm weights without rechecking the first four leaves or the final pole
input would not justify the scalar consumer.

The source checker reconstructs all changed scalar equations, the
unchanged normalized scalar pole equations, the complete point and
subtraction modules, the logarithm circuits, the proxy and its inverse
guard. It also checks the exact affine valuation descent above. Its
finite checks supplement this uniform proof and do not certify the
inherited global arithmetic theorem.
